\documentclass[11pt,a4paper]{article}
\usepackage[utf8]{inputenc}
\usepackage[T1]{fontenc}
\usepackage{geometry}
\usepackage{xcolor}
\usepackage{hyperref}
\usepackage{titlesec}

\geometry{margin=1in}
\definecolor{bosch_red}{RGB}{224, 0, 0}

\hypersetup{
    colorlinks=true,
    linkcolor=bosch_red,
    urlcolor=bosch_red,
}

\titlelabel{\thetitle.\quad}

\begin{document}

\begin{center}
    {\LARGE\bfseries\color{bosch_red} Research Statement Outline}\\[6pt]
    {\large GPAI-Driven Robotic World Models: Bridging System 1/2 for Smart Manufacturing}\\[4pt]
    \textbf{Ismail AISSAOUI} $|$ State Engineer in Artificial Intelligence
\end{center}

\bigskip

\section{Introduction and Vision}
The future of Smart Manufacturing lies in robots that do not just follow static scripts but "reason" about their environment. My research vision focuses on the implementation of **General Purpose AI (GPAI)** that bridges the gap between intuitive perception (**System 1**) and analytical planning (**System 2**). By leveraging **State Space Models (SSMs)** and **Physics-Informed ML**, I aim to develop robotic "World Models" that can autonomously navigate the uncertainty of machine loading and human-robot collaboration in dynamic factory environments.

\section{Proposed Research Axes}

\subsection{1. Hybrid Perception-Planning Architectures (System 1/2)}
I propose to investigate the integration of fast, reactive models (e.g., TCNs or Vision Transformers) for real-time scene understanding with slow, deliberative models (e.g., **Mamba-SSM** or **xLSTM**) for sequential decision-making. This dual-system approach allows the robot to handle high-frequency sensor data while maintaining a long-term "memory" of the manufacturing process, enabling flexible navigation and manipulation strategies that adapt to changing product types.

\subsection{2. Physics-Informed World Models for Robust Manipulation}
Data-driven robotics often struggle with physical consistency. I propose to develop world models that incorporate the **kinematics and physical constraints** of industrial robots directly into the neural architecture using **Physics-Informed Neural Networks (PINN)**. By penalizing violations of physical laws during the "What-If" simulation phase, we can ensure that the GPAI-guided trajectories are not only efficient but also physically feasible and safe, reducing the gap between simulation (Webots/ROS) and reality.

\subsection{3. Explainable and Trustworthy Human-Centric AI}
For robots to work alongside humans, their decisions must be interpretable. I propose to extend my work on **Gradient Saliency** and **Uncertainty Quantification** to the multimodal perception systems of (humanoid) robots. By providing real-time "Confidence Envelopes" and visual heatmaps of the robot's focus, we can create a transparent feedback loop for human operators, ensuring that the GPAI-driven navigation is both trustworthy and ethically aligned with safety standards in a shared manufacturing space.

\section{Alignment with Bosch and HORIZON}
My profile as an AI Engineer with a successful track record in **Hardware-Software Co-Design** and **Industrial Digital Twins** (Sonatrach) aligns perfectly with the requirements of Robert Bosch Manufacturing Solutions. My experience in achieving **0.04ms inference** for edge-based sequential models provides the technical foundation needed to implement these GPAI tools in the **ARENA2036** high-speed manufacturing environment. I am eager to contribute to the HORIZON project’s mission of defining the future of European Smart Manufacturing.

\end{document}
